% =============================================================================
% Appendix subsection — worked example for Figure 2 (running calculator).
% This file is meant to be \input{} into the existing
% `\section{Algorithmic Details}` (label: app:algo) of 9_appendix.tex,
% directly after `app:algo:inferability` so that the worked numbers follow
% the component formulas they instantiate.
% =============================================================================

\subsection{Worked Example: Scoring \texttt{Calculator.total}}
\label{app:algo:worked_example}

To make the numbers shown in Figure~\ref{fig:pdg_score}(b) fully
reproducible, this subsection walks through every quantity that goes
into $\hat{H}$, $\hat{I}$, and the final FIM score on the running
example. The target function is the body shown verbatim in
Figure~\ref{fig:pdg_score}(a):
%
\begin{center}\small
\begin{minipage}{0.78\linewidth}
\begin{verbatim}
def total(self) -> int:
    """Sum positive integer entries."""
    s, n = 0, 0
    for v in self.history:
        if is_int(v):
            if v > 0:
                s = add(s, v)
                n += 1
    if n == 0: return 0
    return s
\end{verbatim}
\end{minipage}
\end{center}

\paragraph{Lines of code (\textnormal{LoC} = 10).}
We follow Python's \texttt{ast} convention,
$\mathrm{LoC}(v)=\texttt{end\_lineno}-(\texttt{lineno}-1)$, i.e.\ the
number of source lines from the \texttt{def} keyword through the last
statement of the body, both inclusive. Counting the ten lines above
verbatim --- including the docstring --- gives $\mathrm{LoC}=10$.

\paragraph{Cyclomatic complexity (\textnormal{CC} = 5).}
The McCabe complexity counter starts at $1$ and adds $1$ per branching
node visited by \texttt{ast.walk}: any \texttt{If}/\texttt{IfExp},
\texttt{For}/\texttt{While}, \texttt{ExceptHandler},
$(\text{values}{-}1)$ per \texttt{BoolOp}, and one per generator inside
a comprehension.
For \texttt{Calculator.total} the contributions are:
%
\begin{center}\small
\begin{tabular}{l l c}
\toprule
Source line & AST node & $\Delta\mathrm{CC}$ \\
\midrule
\textit{(base)}                       & ---                 & $+1$ \\
\texttt{for v in self.history:}       & \texttt{ast.For}    & $+1$ \\
\texttt{if is\_int(v):}               & \texttt{ast.If}     & $+1$ \\
\texttt{if v > 0:}                    & \texttt{ast.If}     & $+1$ \\
\texttt{if n == 0: return 0}          & \texttt{ast.If}     & $+1$ \\
\bottomrule
\end{tabular}
\end{center}
%
yielding $\mathrm{CC}=5$. Tuple unpacking and the augmented assignment
\texttt{n += 1} are not branches and contribute nothing.

\paragraph{Maximum nesting depth (\textnormal{D} = 3).}
The depth counter only increments when entering a structural container
node: \texttt{If}, \texttt{For}, \texttt{While}, \texttt{AsyncFor},
\texttt{With}, \texttt{AsyncWith}, or \texttt{Try}. Crucially, this
list \emph{excludes} \texttt{IfExp} (the inline ``\texttt{a if b else
c}'' expression), so ternary operators do not count. The deepest path
in \texttt{total} is
\texttt{for}~$\to$~\texttt{if is\_int}~$\to$~\texttt{if v>0}, giving
$\mathrm{D}=3$. The trailing \texttt{if n == 0: return 0} sits at the
top level of the function body and reaches depth $1$ only.

\paragraph{Complexity score $\hat{H}=0.40$.}
Plugging into Eq.~\ref{eq:complexity} with the default
$(w_{\ell},w_{c},w_{d})=(0.4,0.4,0.2)$,
$(c_{\ell},c_{c},c_{d})=(50,10,5)$, and $\phi(x,c)=\min(x/c,2)$:
%
\begin{align*}
w_{\ell}\,\phi(10,50)
   &= 0.4\!\cdot\!\min(10/50,2) = 0.4\!\cdot\!0.20 = 0.08, \\
w_{c}\,\phi(5,10)
   &= 0.4\!\cdot\!\min(5/10,2)  = 0.4\!\cdot\!0.50 = 0.20, \\
w_{d}\,\phi(3,5)
   &= 0.2\!\cdot\!\min(3/5,2)   = 0.2\!\cdot\!0.60 = 0.12, \\
\hat{H}(\texttt{total})
   &= 0.08+0.20+0.12 \;=\; 0.40.
\end{align*}

\paragraph{Inferability score $\hat{I}=0.48$.}
We compute each of the five components in Eq.~\ref{eq:inferability}
with the default
$(\alpha,\beta,\gamma,\delta,\varepsilon)=(0.30,0.25,0.20,0.10,0.15)$.

\emph{(i) Caller specificity ($C_{\mathrm{caller}}$).} The only
in-file caller is \texttt{Calculator.mean}, which invokes
\texttt{self.total()} with no arguments and no keywords. The
specificity routine assigns $0.50$ for the call site, $0$ for
arguments and keywords, and a $+0.10$ ``found'' bonus, yielding
$\mathrm{sp}=0.60$. Aggregating over callers and normalizing by $3$,
%
\[
C_{\mathrm{caller}}=\min(0.60/3,\,1.0)=0.20,
\quad
\alpha\,C_{\mathrm{caller}}=0.30\!\cdot\!0.20=0.06.
\]

\emph{(ii) Callee count ($C_{\mathrm{callee}}$).} Inside the body,
\texttt{total} calls \texttt{is\_int} and \texttt{add}---both defined
in the same file---so $|\text{intra-file callees}|=2$. The call
\texttt{self.history.append(v)} does not count because \texttt{append}
is not a member of $\mathcal{V}$. Hence
%
\[
C_{\mathrm{callee}}=\min(2/4,\,1.0)=0.50,
\quad
\beta\,C_{\mathrm{callee}}=0.25\!\cdot\!0.50=0.125\approx 0.13.
\]

\emph{(iii) Signature ($C_{\mathrm{sig}}$).} The four sub-bonuses are
$+0.30$ for the return-type annotation \texttt{-> int}, $0$ for
parameter type annotations (the only argument is \texttt{self}),
$+\min(\text{name parts}/5,0.25)=+0.20$ for the single-token
descriptive name \texttt{total}, and $0$ for non-\texttt{self}
parameters. Thus
%
\[
C_{\mathrm{sig}}=0.30+0.20=0.50,
\quad
\gamma\,C_{\mathrm{sig}}=0.20\!\cdot\!0.50=0.10.
\]

\emph{(iv) Documentation ($C_{\mathrm{doc}}$).} A docstring is
present (the first body statement is a string constant), so
$C_{\mathrm{doc}}=0.5$ and
$\delta\,C_{\mathrm{doc}}=0.10\!\cdot\!0.5=0.05$.

\emph{(v) Class context ($C_{\mathrm{class}}$).} \texttt{total} sits
inside \texttt{Calculator} alongside three sibling methods
(\texttt{\_\_init\_\_}, \texttt{push}, \texttt{mean}). The base score
is $\min(3/5,1.0)=0.60$; because \texttt{Calculator.\_\_init\_\_}
exists and is distinct from \texttt{total}, a $+0.30$ bonus brings
$C_{\mathrm{class}}=0.90$, capped at $1.0$. So
%
\[
\varepsilon\,C_{\mathrm{class}}=0.15\!\cdot\!0.90=0.135\approx 0.14.
\]

Summing the five contributions reproduces the lower bar of
Figure~\ref{fig:pdg_score}(b):
%
\[
\hat{I}(\texttt{total}) \;=\; 0.06+0.13+0.10+0.05+0.14 \;=\; 0.48.
\]
%
The class-context term dominates because \texttt{total} is embedded in
a four-method class that also exposes a constructor; the docstring
term is the smallest because the docstring component itself is capped
at $0.5$ before being multiplied by the small weight $\delta=0.10$.

\paragraph{Final FIM score and selection.}
With $\Delta(\texttt{total})=\max(0,\hat{H}-\hat{I})/(\hat{H}+\epsilon)
=\max(0,-0.08)/0.40=0$, the difficulty penalty
$\rho(\Delta)=1$ and the harmonic-mean-like product gives
%
\[
\mathrm{FIM}(\texttt{total})
\;=\;\frac{0.40\!\times\!0.48}{0.40+0.48+\epsilon}
\;\approx\;\frac{0.192}{0.88}
\;\approx\;0.22
\;\geq\;\tau\!=\!0.20,
\]
%
so \texttt{Calculator.total} is selected as a mask target. The four
short helpers in the file (\texttt{add}, \texttt{is\_int},
\texttt{Calculator.push}, \texttt{Calculator.mean}) are removed before
scoring by the LoC hard filter ($\mathrm{LoC}<10$); the dunder method
\texttt{Calculator.\_\_init\_\_} is removed by the
\texttt{SKIP\_NAMES} list. The geometry of the threshold curve in
Figure~\ref{fig:pdg_score}(c) explains why a moderate-on-both-axes
function like \texttt{total} clears $\tau$ even though no individual
score is large: the harmonic-mean form rewards balanced contributions
rather than a single tall component.
